\VerdierPhysicalPageAnchor{0041}%
\VerdierSemanticAnchor{ch1-1-3-6-eq2}
tels que les diagrammes ci-après soient commutatifs :

\begin{equation}
  \begin{tikzcd}[column sep=7em,row sep=3em]
    FT(hg) \arrow[r,"m(hg)"] \arrow[d,equal] &
      T'(hg)F \arrow[d,equal] \\
    FT(g)T(h) \arrow[d,"m(g)\mathbin{\star}T(h)"'] &
      T'(g)T'(h)F \\
    T'(g)FT(h) \arrow[ur,"T'(g)\mathbin{\star}m(h)"'] &
  \end{tikzcd}
  \qquad g,h\in G.
  \tag{1.3.6.2}
\end{equation}

\VerdierSemanticAnchor{ch1-1-3-6-remark}
Il convient de remarquer que, même dans le cas d'opération stricte de
$G^\circ$ sur les catégories, les foncteurs ne commutent pas
nécessairement strictement aux opérations de $G^\circ$. Ainsi, lorsque
$G^\circ$ opère trivialement sur les catégories $\mathcal C$ et
$\mathcal C'$ (\textit{i.e.} $T(g)=T'(g)=$ identité), un
$G^\circ$-foncteur n'est alors autre qu'un foncteur muni d'une
représentation de $G$ dans son groupe des automorphismes. De plus, il est
facile de voir qu'un $G^\circ$-foncteur n'est pas nécessairement
isomorphe à un $G^\circ$-foncteur strict. On construit facilement un
contre-exemple en s'inspirant de [19].

\VerdierSemanticAnchor{ch1-1-3-6-strictify}
En utilisant les résultats de [\emph{loc. cit.}], on peut d'ailleurs
montrer que toute $G^\circ$-catégorie est $G^\circ$-équivalente à une
$G^\circ$-catégorie stricte. On peut aussi par les mêmes méthodes
préciser ce résultat de la façon suivante. Étant données deux
$G^\circ$-catégories $(\mathcal C,G^\circ,T,c)$ et
$(\mathcal C',G^\circ,T',c')$, il existe deux $G^\circ$-catégories
strictes : $(\widetilde{\mathcal C},G^\circ,\widetilde T)$ et
$(\widetilde{\mathcal C}',G^\circ,\widetilde T')$ et des
$G^\circ$-équivalences :

\VerdierSemanticAnchor{ch1-1-3-6-equivalences}
\[
  (u,\mu)\colon
  (\widetilde{\mathcal C},G^\circ,\widetilde T)
  \longrightarrow
  (\mathcal C,G^\circ,T,c),
\]
\[
  (v,\nu)\colon
  (\mathcal C',G^\circ,T',c')
  \longrightarrow
  (\widetilde{\mathcal C}',G^\circ,\widetilde T')
\]
telles que pour tout foncteur :

\VerdierSemanticAnchor{ch1-1-3-6-functor}
\[
  (F,m)\colon
  (\mathcal C,G^\circ,T,c)
  \longrightarrow
  (\mathcal C',G^\circ,T',c'),
\]
le foncteur composé :

\VerdierSemanticAnchor{ch1-1-3-6-composite}
\[
  (v,\nu)(F,m)(u,\mu)\colon
  (\widetilde{\mathcal C},G^\circ,\widetilde T)
  \longrightarrow
  (\widetilde{\mathcal C}',G^\circ,\widetilde T')
\]
soit isomorphe à un foncteur commutant strictement aux opérations de
$G^\circ$. Nous n'utiliserons pas ce résultat.

\VerdierSemanticAnchor{ch1-1-3-7}
\noindent\textbf{1.3.7.}\quad
Nous avons rappelé dans ce numéro l'équivalence de deux langages : le
langage des catégories graduées et celui des $G^\circ$-catégories. Nous
emploierons
